% Figure: a plane electromagnetic wave. E in the x-z plane, B in the y-z
% plane, propagation along z, E perpendicular to B perpendicular to k.
\begin{figure}[htbp]
\centering
\begin{tikzpicture}
\begin{axis}[
  width=13cm, height=6.5cm,
  xlabel={propagation direction $z$},
  axis lines=middle,
  xmin=0, xmax=12.6, ymin=-1.4, ymax=1.4,
  xtick=\empty, ytick=\empty,
  clip=false,
  label style={font=\small},
]
% E field: sine in the vertical plane
\addplot[engred, very thick, samples=160, domain=0:12.566] {sin(deg(x))};
\node[engred, font=\small] at (axis cs:1.7,1.25) {$\vect{E}$};
% B field: cosine? in phase with E for a wave, drawn as a phase-locked sine
% scaled to look like the transverse companion (shown smaller for the page)
\addplot[engblue, very thick, samples=160, domain=0:12.566] {0.55*sin(deg(x))};
\node[engblue, font=\small] at (axis cs:1.7,0.78) {$\vect{B}$};
% propagation arrow
\draw[->, thick] (axis cs:11.0,0) -- (axis cs:12.3,0) node[above, font=\small] {$c$};
\end{axis}
\end{tikzpicture}
\caption[Plane electromagnetic wave]{A plane electromagnetic wave in
vacuum. The electric field $\vect{E}$ and magnetic field $\vect{B}$
oscillate in phase, transverse to each other and to the propagation
direction, with $|\vect{E}|/|\vect{B}| = c = 1/\sqrt{\mu_0\varepsilon_0}$.
The two curl equations of Maxwell, Faraday and Amp\`ere--Maxwell, feed
each other: a changing $\vect{E}$ sources $\vect{B}$ and a changing
$\vect{B}$ sources $\vect{E}$, and the pair propagates at $c$. The energy
density splits equally between the two fields, and the time-averaged
Poynting vector $\langle S\rangle = \tfrac12\varepsilon_0 c E_0^2$ is the
intensity the wave delivers. The two curves are drawn at different
heights only for legibility; their amplitudes are fixed by the ratio
$c$.}
\label{fig:vol04ch08:em-wave}
\end{figure}
